\section{Related Work}
% \subsection{Intent Classification}
% Intent classification is a typical text classification task where class labels are intent names. Various neural networks have inspired a wide range of studies aimed at creating neural models for text classification. Various neural model structures such as CNN \cite{conneau2016very,kim-2014-convolutional}, LSTM \cite{yang2016hierarchical} and GCN \cite{yao2019graph,liu2024spatial} have proven more effective than conventional methods. Additionally, some studies \cite{zhang2017multi,zhang2020multi} utilize label embeddings and train them simultaneously with the input texts. In recent developments, the accomplishments of large-scale pre-training language models (PLM) \cite{kenton2019bert} have generated significant interest in incorporating this pre-training approach \cite{reimers2019sentence} into monolingual and multilingual text classification \cite{conneau2020unsupervised}. However, most works are single-turn text classification tasks, which is unsuitable for multilingual multi-turn intent classification.

\textbf{Modeling Multi-turn Dialogue Context}: Modelling the multi-turn dialogues is the foundation for dialogue understanding tasks. Previous works adopt bidirectional contextual LSTM \cite{ghosal2021exploring,liu2022title2vec} to create context-aware utterance representation on MultiWOZ intent classification \cite{budzianowski2018multiwoz}. Recent works use PLM as a sentence encoder \cite{shen2021directed} on emotion recognition in conversation (ERC). Specifically, \cite{lee2022compm} used PLM to encode the context and speaker's memory and \cite{qin2023bert} enhance PLM by integrating multi-turn info from the utterance, context and dialogue structure through fine-tuning. However, all of their tasks adopt the multi-turn dialogue training set, which is hard to collect for an e-commerce chatbot. Our method attempts to combine an XLM-based model trained on the single-turn dataset into an in-context retrieval augmented pipeline with LLM, solving the multi-turn intent classification task in a zero-shot setting. 

%\subsection{In-context learning and Retrieval Augmented Generation}
% will link the references later
%In-context learning has become successful with the Large language model and has changed the approach to text classification. Not like fine-tuning, struggles when the number of available training samples is limited \cite{dodge2020fine}. The scarcity of contextual multi-turn dialogue interaction further calls for the need of methods which do not rely on expensive training data curation. \cite{brown2020language} has shown that large-scale language models performs competitively in many NLP tasks through in-context learning (ICL), even when no explicit training on the task was done. Recent work \cite{milios2023context} has applied ICL to a classification task with many labels. In this approach, demonstrations are dynamically retrieved to fit within the model's context window. Our work diverges from this in that the demonstrations in the context are not precisely aligned with the target task; while the demonstrations focus on single-turn intent recognition, the target task involves multi-turn intent recognition. However, \cite{cheng2023uprise} has also demonstrated that even when the demonstrations do not precisely match the target task, they can still aid in model inference, especially if they are closely related.

\noindent\textbf{In-context Retrieval}: In-context learning (ICL) with LLM like GPT-3 \cite{gpt3} demonstrates the significant improvement on few-shot/zero-shot NLP tasks. ICL has been successful in utterance-level tasks like intent classification \cite{Yu2021FewshotIC}. As for the \textbf{Retrieval} part, most research on in-context learning (ICL) usually deals with single sentences or document retrieval, but we are interested in finding and understanding dialogues. Generally, there are two types of systems to find the relevant dialogues: the first is LM-score based retrieval. They \cite{Rubin2021LearningTR, Shin2021ConstrainedLM} check the probability of a language model, like GPT-3, to decode the right answer based on an example. The second type defines similarity metrics between task results and uses them as the training objective for the retriever. Both K-highest and lowest examples are used as positive and negative samples to help the system learn. The most pertinent research on dialogue retrieval concentrates on areas such as knowledge identification \cite{Wu2021DIALKIKI} and response selection \cite{Yuan2019MultihopSN}. Our objectives and settings differ from them.